% ЗАКЛЮЧЕНИЕ, без номер.
% Обобщава решението, предимствата и недостатъците му, резултатите и посоките напред.
% Обобщението на резултатите се дописва СЛЕД като Раздел VI бъде попълнен.
\section*{ЗАКЛЮЧЕНИЕ}
\label{sec:conclusion}
\addcontentsline{toc}{section}{ЗАКЛЮЧЕНИЕ}

Проектът разглежда конвейер за обработка на облаци от точки и карти на дълбочината, в
който търсенето на съседи е обособено като отделен компонент върху пространствен индекс
и е ускорено чрез явна векторизация. Обособяването има две последици. Първата е
инженерна: всяка стъпка от конвейера ползва един и същ тесен интерфейс, така че
реализацията на индекса се сменя без промени в стъпките. Втората е методическа: цената
на всяка стъпка се измерва самостоятелно, а двете реализации на индекса се сравняват
върху едни и същи данни.

Основното ограничение на решението е, че печалбата от векторизацията зависи от
разположението на данните по компонента, а то прави достъпа до отделна точка по-скъп.
Изборът е изгоден за конвейер, в който преобладават заявките по околност, и неизгоден
за натоварване, в което се четат единични точки в произволен ред. Второ ограничение е
съзнателното изключване на разпознаването по обучени модели, което оставя последното
ниво на веригата извън обхвата и е изложено в увода.

\TODO{обобщение на измерените резултати. Пише се СЛЕД попълването на Раздел
\textbackslash ref\{sec:results\} и не съдържа твърдение, което не се вижда в таблица оттам}

Посоките за продължение са три. Първата е приблизително търсене на съседи с контролирана
загуба на точност~\cite{muja2014flann}, което е разгледано и оставено извън обхвата
именно защото внася втори източник на грешка. Втората е паралелно изпълнение на
стъпките върху няколко ядра, за което текущият договор на стъпката вече е подходящ,
понеже стъпките не променят входа си на място. Третата е добавяне на разпознавател над
изхода на конвейера, което би затворило веригата от сензор до решение.
